\chapter{Gaussian Elimination and Linear Systems}

Every later chapter of this volume quietly assumes this one. The inverse of
a $3\times3$ matrix is computed by row-reducing $[A\mid I]$; the rank of a
matrix is the number of non-zero rows after row reduction; the simplex
method of the linear-programming volume is row reduction with a rule for
choosing which pivot to use next. It is worth having the mechanism itself
drilled until it is automatic.

The idea is older and simpler than matrices. To solve a system of linear
equations one adds multiples of equations to each other until each new
equation contains one unknown fewer, then reads the unknowns off backwards.
Gaussian elimination is that procedure written down without the letters:
only the coefficients are carried along, in a rectangular array, and the
operations on equations become operations on \emph{rows}.

Two things make it more than a computational recipe. It never loses or
invents solutions --- each step is reversible, so the solution set is
untouched --- and the shape of the final array reveals immediately which of
the three possible outcomes holds: one solution, none, or infinitely many.
That last case is where the parameters come from, and it is the case
examinations ask about.

\section{Core Rules}

\begin{rulebox}{The augmented matrix}
The system
\[
  \begin{array}{r@{}c@{}r@{}c@{}r@{}c@{}l}
    a_{11}x_{1}&{}+{}&\cdots&{}+{}&a_{1n}x_{n}&{}={}&b_{1}\\
    &&\vdots&&&&\\
    a_{m1}x_{1}&{}+{}&\cdots&{}+{}&a_{mn}x_{n}&{}={}&b_{m}
  \end{array}
\]
is recorded as the \textbf{augmented matrix} $[A\mid \mathbf{b}]$: the
coefficients on the left of the bar, the right-hand sides on the right. The
unknowns are not written; their identity is their column.
\end{rulebox}

\begin{rulebox}{The three elementary row operations}
\begin{enumerate}[leftmargin=1.8em,itemsep=0.3ex]
  \item \textbf{Swap} two rows. \hfill $R_{i}\leftrightarrow R_{j}$
  \item \textbf{Scale} a row by a constant $c\neq0$. \hfill $R_{i}\to cR_{i}$
  \item \textbf{Add} a multiple of one row to another. \hfill
        $R_{i}\to R_{i}+cR_{j}$
\end{enumerate}
Each corresponds to something obviously harmless done to the equations ---
reordering them, multiplying one by a non-zero number, adding one to
another --- and each is \textbf{reversible}. Hence
\[
  [A\mid\mathbf{b}]\ \longrightarrow\ [A'\mid\mathbf{b}']
  \quad\Longrightarrow\quad
  \text{same solution set.}
\]
The requirement $c\neq0$ in the second operation is what keeps it
reversible: multiplying a row by $0$ destroys information.
\end{rulebox}

\begin{rulebox}{Row-echelon form and pivots}
A matrix is in \textbf{row-echelon form} (REF) if
\begin{itemize}[leftmargin=1.6em,itemsep=0.2ex]
  \item all zero rows are at the bottom, and
  \item in every non-zero row the first non-zero entry --- its
        \textbf{pivot} --- stands strictly to the right of the pivot of the
        row above.
\end{itemize}
The result is a staircase descending to the right. In \textbf{reduced}
row-echelon form (RREF) additionally every pivot equals $1$ and is the only
non-zero entry in its column.
\[
  \underbrace{\begin{pmatrix}
    \mathbf{2} & 1 & -1 & 3\\
    0 & \mathbf{3} & 4 & 1\\
    0 & 0 & 0 & 0
  \end{pmatrix}}_{\text{REF, pivots }2,3}
  \qquad
  \underbrace{\begin{pmatrix}
    \mathbf{1} & 0 & 5 & 2\\
    0 & \mathbf{1} & 4 & 1\\
    0 & 0 & 0 & 0
  \end{pmatrix}}_{\text{RREF}}
\]
A column containing a pivot is a \textbf{pivot column}; its variable is
determined. A non-pivot column belongs to a \textbf{free variable}, which
may be chosen at will.
\end{rulebox}

\begin{rulebox}{The algorithm}
\textbf{Forward elimination.} Work column by column, left to right.
\begin{enumerate}[leftmargin=1.8em,itemsep=0.2ex]
  \item Find the topmost unused row with a non-zero entry in the current
        column; swap it into position. (If the whole column is zero, move
        to the next column.)
  \item Use operation 3 to clear every entry \emph{below} that pivot.
  \item Move down one row and right one column, and repeat.
\end{enumerate}
\textbf{Back substitution.} From the bottom pivot upwards, solve each row
for its pivot variable and substitute into the rows above. Alternatively
continue with row operations until the matrix is in RREF --- then each
pivot row reads off its variable directly. That variant is
\textbf{Gauss--Jordan}.

Arithmetic advice that pays for itself: keep integers as long as possible
by scaling rows rather than dividing early, and never clear a column using
a row you have not finished with.
\end{rulebox}

\begin{rulebox}{Reading the solution set off the echelon form}
Let $r$ be the number of pivots in $A$ (that is, $\mathrm{rk}(A)$) and $n$
the number of unknowns.
\begin{itemize}[leftmargin=1.6em,itemsep=0.35ex]
  \item \textbf{No solution} exactly when a row reads
        $\bigl(0\ \cdots\ 0 \mid \beta\bigr)$ with $\beta\neq0$: the
        equation $0=\beta$ is false. Equivalently
        $\mathrm{rk}([A\mid\mathbf{b}])>\mathrm{rk}(A)$. The system is
        \emph{inconsistent}.
  \item \textbf{Exactly one solution} when the system is consistent and
        $r=n$: every variable is a pivot variable.
  \item \textbf{Infinitely many} when the system is consistent and $r<n$:
        there are $n-r$ free variables. Setting them to parameters
        $t_{1},\dots,t_{n-r}$ gives the general solution as one particular
        solution plus a combination of $n-r$ direction vectors.
\end{itemize}
A \textbf{homogeneous} system ($\mathbf{b}=\mathbf{0}$) is never
inconsistent --- $\mathbf{x}=\mathbf{0}$ always solves it --- so it has
non-trivial solutions exactly when $r<n$.
\end{rulebox}

\begin{rulebox}{Two by-products}
\begin{itemize}[leftmargin=1.6em,itemsep=0.35ex]
  \item \textbf{The inverse.} Row-reduce $[A\mid I]$ to RREF. If the left
        block becomes $I$, the right block is $A^{-1}$; if a zero row
        appears on the left, $A$ is singular and has no inverse. This is the
        practical method for $n\ge3$, where the adjugate formula becomes
        unwieldy.
  \item \textbf{The determinant.} Row operations change $\det$ in a known
        way: a swap flips its sign, scaling a row by $c$ multiplies it by
        $c$, and adding a multiple of one row to another leaves it
        unchanged. So reducing to a triangular form and multiplying the
        pivots --- correcting for swaps and scalings --- computes the
        determinant in far fewer operations than cofactor expansion.
\end{itemize}
\end{rulebox}

\section{Worked Example}

\begin{worked}{Solve $x+y+z=6$, $2x-y+z=3$, $x+2y-z=2$ by elimination --- and then change one number so that the system breaks}
\wstep{Step 1. Write the augmented matrix. Never carry the letters.}
\[
  \begin{array}{rcl}
    x+y+z&=&6\\ 2x-y+z&=&3\\ x+2y-z&=&2
  \end{array}
  \qquad\Longrightarrow\qquad
  [A\mid\mathbf b]=
  \left(\begin{array}{ccc|c}
    1&1&1&6\\ 2&-1&1&3\\ 1&2&-1&2
  \end{array}\right)
\]
The vertical bar is not decoration: everything left of it is a coefficient,
everything right of it is a constant, and row operations act on the whole
row including the constant.
\wtrap{Line the variables up in the same order in every equation, and write
a $0$ where a variable is absent. A missing zero shifts a whole column and
corrupts everything downstream.}

\wstep{Step 2. Clear the first column using the first row.}
$R_{2}\to R_{2}-2R_{1}$ and $R_{3}\to R_{3}-R_{1}$:
\[
  \left(\begin{array}{ccc|c}
    1&1&1&6\\ 0&-3&-1&-9\\ 0&1&-2&-4
  \end{array}\right)
\]
Check one entry by hand: the new $(2,4)$ entry is $3-2\cdot6=-9$. \checkmark

\wstep{Step 3. Clear the second column using the second row.}
$R_{3}\to 3R_{3}+R_{2}$ avoids fractions:
\[
  \left(\begin{array}{ccc|c}
    1&1&1&6\\ 0&-3&-1&-9\\ 0&0&-7&-21
  \end{array}\right)
\]
\wtrap{Scaling $R_{3}$ by $3$ is harmless \emph{for solving} --- it does not
change the solution set. It does change the determinant, by a factor $3$. If
the same reduction is being used to compute $\det A$, record every scaling
and every swap; if you are only solving, ignore them.}

\wstep{Step 4. Count pivots before back-substituting.}
Three pivots ($1$, $-3$, $-7$) in three columns, and no row of the form
$(0\ 0\ 0\mid c)$ with $c\neq0$. So: consistent, no free variables,
\textbf{exactly one solution}. This diagnosis comes before the arithmetic,
and it is what an examination question usually asks for.

\wstep{Step 5. Back-substitute from the bottom.}
\[
  -7z=-21\;\Rightarrow\;z=3;\qquad
  -3y-3=-9\;\Rightarrow\;y=2;\qquad
  x+2+3=6\;\Rightarrow\;x=1 .
\]
\[
  \boxed{(x,y,z)=(1,2,3)}
\]

\wstep{Step 6. Check in the \emph{original} equations, not the reduced ones.}
$1+2+3=6$ \checkmark;\; $2-2+3=3$ \checkmark;\; $1+4-3=2$ \checkmark.
Checking against the row-echelon form would only confirm the back
substitution, not the elimination --- so it catches half the possible
errors.

\wstep{Step 7. Now break it, to see the other two outcomes.}
Replace the third equation by $3x+3y+3z=18$, a copy of the first. Then
$R_{3}\to R_{3}-3R_{1}$ gives the row $(0\ 0\ 0\mid0)$: two pivots, three
unknowns, so one \textbf{free variable} and infinitely many solutions.
Setting $z=t$ and back-substituting through $-3y-t=-9$ gives
\[
  y=3-\tfrac{t}{3},\qquad x=6-y-z=3-\tfrac{2t}{3},\qquad t\in\R .
\]
Replace it instead by $3x+3y+3z=17$ and the same operation gives
$(0\ 0\ 0\mid-1)$: the equation $0=-1$, so \textbf{no solution}.
\wtrap{$(0\ 0\ 0\mid 0)$ and $(0\ 0\ 0\mid c\neq0)$ look almost identical
and mean opposite things --- infinitely many solutions versus none. Read the
entry to the right of the bar before concluding anything, and never
``simplify away'' a zero row without looking at it.}

\wstep{Step 8. What else the reduction has told you.}
The same triangular form gives $\det A=1\cdot(-3)\cdot(-7)/3=7$ --- dividing
by the $3$ introduced in Step~3 --- and shows $\operatorname{rank}A=3$. One
elimination answers three questions: solve, rank, determinant.
\end{worked}

\section{Practice Problems}

\begin{problems}
  \item Name the three elementary row operations and say why each one
        leaves the solution set unchanged.
  \item Why must the scaling factor in operation 2 be non-zero?
  \item Is $\begin{pmatrix}1&2&3\\0&0&4\\0&5&6\end{pmatrix}$ in row-echelon
        form? If not, make it so with one operation.
  \item Is $\begin{pmatrix}1&0&2\\0&1&3\\0&0&0\end{pmatrix}$ in REF, in
        RREF, or neither?
  \item Give the pivots and the pivot columns of
        $\begin{pmatrix}0&2&1&3\\0&0&0&5\\0&0&0&0\end{pmatrix}$.
  \item How many free variables does a consistent system with $5$ unknowns
        and $3$ pivots have?
  \item A row of an augmented matrix reads $(0\ 0\ 0\mid 4)$. What does that
        tell you?
  \item A row of an augmented matrix reads $(0\ 0\ 0\mid 0)$. What does
        \emph{that} tell you?
\end{problems}

Solve the systems in problems 9--20 by elimination, and state which of the
three outcomes occurs.

\begin{problems}[resume]
  \item $x+y=5$, \;$2x-y=1$.
  \item $2x+3y=12$, \;$x-y=1$.
  \item $x+y+z=6$, \;$2x-y+z=3$, \;$x+2y-z=2$.
  \item $x+2y+3z=14$, \;$2x-y+z=3$, \;$3x+y-z=2$.
  \item $x+y=2$, \;$2x+2y=5$.
  \item $x+y+z=3$, \;$2x+2y+2z=6$.
  \item $x+y-z=1$, \;$2x+3y+z=4$, \;$3x+4y=5$.
  \item $x+2y-z=0$, \;$2x-y+3z=0$.
  \item $2x+y-z=1$, \;$x-y+2z=5$, \;$3x+3y-4z=-3$.
  \item $x+2y=4$, \;$3x-y=5$, \;$2x+3y=7$.
  \item $x-2y+z=3$, \;$2x-4y+2z=6$, \;$-x+2y-z=-3$.
  \item $x+y+2z=1$, \;$2x+3y+5z=3$, \;$x+2y+3z=2$.
\end{problems}

In problems 21--26 give the general solution in parametric form.

\begin{problems}[resume]
  \item $x+2y+z=4$, \;$2x+4y+2z=8$.
  \item $x+y+z+w=2$, \;$x-y+z-w=0$.
  \item $x+y-z=1$, \;$2x+3y+z=4$.
  \item $2x-y+z=0$, \;$4x-2y+2z=0$.
  \item $x+2y-3z=0$, \;$2x+5y-8z=0$.
  \item $x_{1}+x_{2}+x_{3}+x_{4}=0$ \;(one equation, four unknowns).
\end{problems}

In problems 27--30 determine the rank by row reduction.

\begin{problems}[resume]
  \item $\begin{pmatrix}1&2&3\\2&4&6\\1&1&1\end{pmatrix}$
  \item $\begin{pmatrix}1&2\\3&4\\5&6\end{pmatrix}$
  \item $\begin{pmatrix}1&1&1\\1&2&3\\1&3&6\end{pmatrix}$
  \item $\begin{pmatrix}2&-1&3\\4&-2&6\\-2&1&-3\end{pmatrix}$
\end{problems}

In problems 31--34 invert the matrix by row-reducing $[A\mid I]$, or show
that it is singular.

\begin{problems}[resume]
  \item $A=\begin{pmatrix}1&2\\3&4\end{pmatrix}$
  \item $A=\begin{pmatrix}1&1&0\\0&1&1\\1&0&1\end{pmatrix}$
  \item $A=\begin{pmatrix}2&0&0\\0&3&0\\0&0&5\end{pmatrix}$
  \item $A=\begin{pmatrix}1&2\\2&4\end{pmatrix}$
\end{problems}

\begin{problems}[resume]
  \item For which $t$ does $x+y=1$, \;$x+ty=2$ have a unique solution, no
        solution, or infinitely many?
  \item For which $t$ does $x+2y=3$, \;$2x+ty=6$ have a unique solution, no
        solution, or infinitely many?
  \item For which $t$ does the system $x+y+z=1$, \;$x+2y+3z=2$,
        \;$x+3y+tz=3$ fail to have a unique solution, and what happens
        then?
  \item For which $t$ does $tx+y=1$, \;$x+ty=1$ have infinitely many
        solutions?
  \item Compute $\det\begin{pmatrix}1&2&3\\2&5&3\\1&0&8\end{pmatrix}$ by
        row-reducing to triangular form, tracking the effect of each
        operation on the determinant.
  \item A $4\times4$ matrix is row-reduced to an upper triangular matrix
        using two row swaps, one scaling by $\tfrac12$, and several
        additions. The resulting diagonal is $(2,3,1,4)$. What was
        $\det A$?
  \item Explain why a homogeneous system with more unknowns than equations
        always has a non-trivial solution.
  \item Show that if $A\mathbf{x}=\mathbf{b}$ has two distinct solutions,
        it has infinitely many.
\end{problems}

\begin{problems}[resume]
  \item \textbf{Leontief.} With input matrix
        $A=\begin{pmatrix}0.2&0.4\\0.3&0.1\end{pmatrix}$ and final demand
        $\mathbf{d}=\binom{12}{9}$, solve
        $(I-A)\mathbf{x}=\mathbf{d}$ by elimination.
  \item \textbf{Three-good equilibrium.} Solve
        $2p_{1}-p_{2}=5$, \;$-p_{1}+3p_{2}-p_{3}=12$,
        \;$-p_{2}+2p_{3}=15$.
  \item \textbf{National income.} With $Y=C+I+G$,
        $C=100+0.8(Y-T)$, $T=0.25Y$, $I=150$, $G=200$: write the model as a
        linear system in $(Y,C,T)$ and solve it by elimination.
\end{problems}

\section{Solutions}

\begin{solutions}
  \item Swapping two rows reorders the equations; scaling a row by $c\neq0$
        multiplies one equation by a non-zero constant; adding a multiple of
        one row to another replaces one equation by a consequence of two of
        them. Every one is reversible --- swap back, scale by $1/c$,
        subtract the same multiple --- so no solution can be lost or
        gained.
  \item Because multiplying by $0$ is not reversible: the row becomes
        $0=0$, and the information in that equation is destroyed. Solutions
        would be gained.
  \item Not in REF: the pivot of row 2 sits in column 3, but row 3 has a
        non-zero entry in column 2, which is further left. Swapping
        $R_{2}\leftrightarrow R_{3}$ gives
        $\begin{pmatrix}1&2&3\\0&5&6\\0&0&4\end{pmatrix}$, which is.
  \item Both: pivots are $1$, each is alone in its column, and the zero row
        is at the bottom. It is in RREF (hence also in REF).
  \item Pivots $2$ (row 1, column 2) and $5$ (row 2, column 4). Pivot
        columns are $2$ and $4$; columns $1$ and $3$ are free.
  \item $n-r=5-3=2$ free variables.
  \item The equation $0=4$: the system is \textbf{inconsistent} and has no
        solution.
  \item Nothing at all --- $0=0$ is true for every $\mathbf{x}$. The row is
        redundant, which signals that the original equations were
        dependent, so a free variable is likely.
  \item $R_{2}\to R_{2}-2R_{1}$ gives $-3y=-9$, $y=3$, then $x=2$. Unique:
        $(x,y)=(2,3)$.
  \item From the second equation $x=y+1$; substituting,
        $2y+2+3y=12$ so $y=2$, $x=3$. Unique.
  \item The worked example: $(x,y,z)=(1,2,3)$. Unique.
  \item Eliminating gives $y+z=5$ and $y+2z=8$, so $z=3$, $y=2$,
        $x=14-4-9=1$. Unique: $(1,2,3)$.
  \item $R_{2}\to R_{2}-2R_{1}$ leaves $(0\ 0\mid1)$, i.e. $0=1$:
        \textbf{no solution}.
  \item $R_{2}\to R_{2}-2R_{1}$ leaves a zero row. One pivot, three
        unknowns: \textbf{infinitely many}, $x=3-y-z$ with $y,z$ free.
  \item $R_{3}=R_{1}+R_{2}$, so the third equation is redundant. Two
        pivots, three unknowns: infinitely many. With $z=t$:
        $y=2-3t$ and $x=-1+4t$, i.e.
        $(x,y,z)=(-1,2,0)+t(4,-3,1)$.
  \item Homogeneous, two equations, three unknowns, so non-trivial
        solutions exist. Eliminating gives $-5y+5z=0$, so $y=z$ and
        $x=z-2y=-z$: the solution set is $t(-1,1,1)$.
  \item $R_{2}\to 2R_{2}-R_{1}$ gives $-3y+5z=9$;
        $R_{3}\to 2R_{3}-3R_{1}$ gives $3y-5z=-9$, the same equation.
        Two pivots, three unknowns: infinitely many. With $z=t$:
        $y=\tfrac{5t-9}{3}$ and, from the first equation,
        $x=\tfrac{1+z-y\cdot 1}{2}$ --- explicitly
        $x=\tfrac{1}{2}\bigl(1+t-\tfrac{5t-9}{3}\bigr)
        =\tfrac{12-2t}{6}=2-\tfrac{t}{3}$. So
        $(x,y,z)=(2,-3,0)+t\bigl(-\tfrac13,\tfrac53,1\bigr)$.
  \item Three equations, two unknowns. From the first two: $x=2$, $y=1$.
        Testing in the third: $4+3=7$ \checkmark, so the third is
        consistent and the unique solution is $(2,1)$. (An
        over-determined system need not be inconsistent.)
  \item The second row is twice the first and the third is minus the first,
        so both reduce to zero rows. One pivot, three unknowns: infinitely
        many, $x=3+2y-z$ with $y,z$ free.
  \item Eliminating: $R_{2}\to R_{2}-2R_{1}$ gives $y+z=1$;
        $R_{3}\to R_{3}-R_{1}$ gives $y+z=1$ --- the same. Two pivots,
        three unknowns: infinitely many. With $z=t$: $y=1-t$ and
        $x=1-y-2t=-t$, so $(x,y,z)=(0,1,0)+t(-1,-1,1)$.
  \item The second equation is twice the first, so one pivot and two free
        variables: $x=4-2s-t$, $y=s$, $z=t$, i.e.
        $(4,0,0)+s(-2,1,0)+t(-1,0,1)$.
  \item Adding the equations gives $2x+2z=2$, so $x+z=1$; subtracting gives
        $2y+2w=2$, so $y+w=1$. With $z=s$ and $w=t$:
        $(x,y,z,w)=(1-s,\,1-t,\,s,\,t)$.
  \item $R_{2}\to R_{2}-2R_{1}$ gives $y+3z=2$, so $y=2-3t$ with $z=t$, and
        $x=1-y+z=1-(2-3t)+t=-1+4t$:
        $(x,y,z)=(-1,2,0)+t(4,-3,1)$.
  \item The second equation is twice the first, so there is one pivot and
        two free variables. Choosing $x$ and $z$ as the free ones avoids
        fractions: $y=2x+z$, so with $x=s$ and $z=t$
        \[
          (x,y,z)=(s,\;2s+t,\;t)=s(1,2,0)+t(0,1,1).
        \]
        The solution set is the plane $2x-y+z=0$ through the origin, and
        any two independent directions in it give an equally valid
        answer.
  \item $R_{2}\to R_{2}-2R_{1}$ gives $y-2z=0$, so $y=2t$ with $z=t$, and
        $x=3z-2y=3t-4t=-t$: $(x,y,z)=t(-1,2,1)$.
  \item One pivot, four unknowns, three free:
        $x_{1}=-s-t-u$, $x_{2}=s$, $x_{3}=t$, $x_{4}=u$.
  \item $R_{2}\to R_{2}-2R_{1}$ gives a zero row;
        $R_{3}\to R_{3}-R_{1}$ gives $(0,-1,-2)$. Two non-zero rows:
        $\mathrm{rk}=2$.
  \item $R_{2}\to R_{2}-3R_{1}$ and $R_{3}\to R_{3}-5R_{1}$ give
        $(0,-2)$ and $(0,-4)$; then $R_{3}\to R_{3}-2R_{2}$ gives a zero
        row. $\mathrm{rk}=2$.
  \item Subtracting $R_{1}$ from the other rows gives $(0,1,2)$ and
        $(0,2,5)$; then $R_{3}\to R_{3}-2R_{2}$ gives $(0,0,1)$. Three
        pivots: $\mathrm{rk}=3$.
  \item $R_{2}=2R_{1}$ and $R_{3}=-R_{1}$, so both become zero rows.
        $\mathrm{rk}=1$.
  \item $\left(\begin{array}{cc|cc}1&2&1&0\\3&4&0&1\end{array}\right)
        \to R_{2}-3R_{1}:\ (0,-2\mid-3,1)
        \to -\tfrac12R_{2}:\ (0,1\mid\tfrac32,-\tfrac12)
        \to R_{1}-2R_{2}$, giving
        $A^{-1}=\begin{pmatrix}-2&1\\ \tfrac32&-\tfrac12\end{pmatrix}$.
  \item Row-reducing $[A\mid I]$ gives
        $A^{-1}=\tfrac12\begin{pmatrix}1&-1&1\\1&1&-1\\-1&1&1\end{pmatrix}$.
        (Check: $A A^{-1}=I$, e.g. row $1$ of $A$ times column $1$ of
        $A^{-1}$ is $\tfrac12(1+1+0)=1$.)
  \item Diagonal, so each row is scaled independently:
        $A^{-1}=\mathrm{diag}\bigl(\tfrac12,\tfrac13,\tfrac15\bigr)$.
  \item $R_{2}\to R_{2}-2R_{1}$ turns the left block's second row into
        $(0,0)$. A zero row on the left means the reduction cannot reach
        $I$: $A$ is \textbf{singular} and has no inverse (indeed
        $\det A=4-4=0$).
  \item Eliminating gives $(t-1)y=1$. If $t\neq1$ there is a unique
        solution, $y=\tfrac{1}{t-1}$ and $x=1-y$. If $t=1$ the row reads
        $0=1$: no solution. Infinitely many never occurs.
  \item Eliminating gives $(t-4)y=0$. If $t\neq4$: unique, $y=0$ and
        $x=3$. If $t=4$ the second equation is exactly twice the first, so
        there are infinitely many solutions, $x=3-2y$ with $y$ free. No
        solution never occurs.
  \item Subtracting $R_{1}$ from $R_{2}$ and $R_{3}$ gives $y+2z=1$ and
        $2y+(t-1)z=2$; then $R_{3}\to R_{3}-2R_{2}$ gives $(t-5)z=0$. For
        $t\neq5$ the solution is unique, $z=0$, $y=1$, $x=0$. For $t=5$ the
        last row vanishes and $z$ is free: infinitely many, namely
        $(x,y,z)=(0,1,0)+z(1,-2,1)$. The system is never inconsistent.
  \item The determinant is $t^{2}-1$, so the system is degenerate for
        $t=\pm1$. For $t=1$ both equations read $x+y=1$: infinitely many.
        For $t=-1$ they read $-x+y=1$ and $x-y=1$, i.e. $0=2$: no
        solution. So infinitely many exactly when $t=1$.
  \item $R_{2}\to R_{2}-2R_{1}$ and $R_{3}\to R_{3}-R_{1}$ (both leave
        $\det$ unchanged) give
        $\begin{pmatrix}1&2&3\\0&1&-3\\0&-2&5\end{pmatrix}$, then
        $R_{3}\to R_{3}+2R_{2}$ gives
        $\begin{pmatrix}1&2&3\\0&1&-3\\0&0&-1\end{pmatrix}$. No swaps and
        no scalings were used, so $\det A=1\cdot1\cdot(-1)=-1$.
  \item Each swap multiplies the determinant by $-1$ (two swaps: no net
        change) and scaling a row by $\tfrac12$ halves it, while the
        additions change nothing. So
        $\tfrac12\det A=2\cdot3\cdot1\cdot4=24$, giving $\det A=48$.
  \item With $m$ equations and $n>m$ unknowns the number of pivots is at
        most $m<n$, so at least $n-m\ge1$ variables are free. A homogeneous
        system is always consistent, so those free variables generate
        infinitely many solutions, all but one of them non-trivial.
  \item Let $\mathbf{x}_{1}\neq\mathbf{x}_{2}$ both solve it and set
        $\mathbf{d}=\mathbf{x}_{1}-\mathbf{x}_{2}\neq\mathbf{0}$. Then
        $A\mathbf{d}=\mathbf{b}-\mathbf{b}=\mathbf{0}$, so for every
        $t\in\R$,
        $A(\mathbf{x}_{1}+t\mathbf{d})=\mathbf{b}+t\mathbf{0}=\mathbf{b}$.
        Every $t$ gives a solution, and they are all distinct.
  \item $I-A=\begin{pmatrix}0.8&-0.4\\-0.3&0.9\end{pmatrix}$. Scaling both
        rows by $10$ and eliminating:
        $8x-4y=120$, $-3x+9y=90$; from the first $y=2x-30$, and then
        $-3x+18x-270=90$ gives $15x=360$, $x=24$ and $y=18$. Gross outputs
        $\mathbf{x}=(24,18)$. (Check: $0.8\cdot24-0.4\cdot18=12$
        \checkmark.)
  \item From the first, $p_{2}=2p_{1}-5$; from the third,
        $p_{3}=\tfrac{15+p_{2}}{2}=5+p_{1}$. Substituting into the second:
        $-p_{1}+3(2p_{1}-5)-(5+p_{1})=4p_{1}-20=12$, so $p_{1}=8$,
        $p_{2}=11$, $p_{3}=13$. (Check the third: $-11+26=15$
        \checkmark.)
  \item Written as a system:
        \[
          \begin{array}{rcl}
            Y-C&=&350\\
            C-0.8Y+0.8T&=&100\\
            T-0.25Y&=&0
          \end{array}
        \]
        Substituting $T=0.25Y$ into the second gives
        $C=100+0.8Y-0.2Y=100+0.6Y$; the first then gives
        $Y-100-0.6Y=350$, so $0.4Y=450$ and $Y=1125$. Hence
        $T=281.25$ and $C=775$. (Check: $775+150+200=1125$ \checkmark.)
        The multiplier $1/0.4=2.5$ is the coefficient sitting in that
        elimination step.
\end{solutions}

\section{Additional Problems}

\begin{problems}
  \item Bring $\begin{pmatrix}0&1&2\\1&2&3\\2&4&6\end{pmatrix}$ to REF and
        give its rank.
  \item Bring the matrix of the previous problem to RREF.
  \item Solve $3x+2y=7$, \;$x-y=-1$.
  \item Solve $x+y+z=1$, \;$x-y+2z=4$, \;$2x+y-z=-1$.
  \item Solve $2x+4y-2z=2$, \;$x+2y-z=1$, \;$3x+6y-3z=3$.
  \item Solve $x+y=1$, \;$y+z=2$, \;$x+z=3$.
  \item Solve $x+2y+3z=1$, \;$4x+5y+6z=2$, \;$7x+8y+9z=4$.
  \item Solve the homogeneous system $x+2y+3z=0$, \;$2x+4y+6z=0$.
  \item Solve $x_{1}-x_{2}+2x_{3}=1$, \;$2x_{1}+x_{3}=3$.
  \item Give the general solution of $x+y+z=6$, \;$x-y+z=2$ in parametric
        form.
  \item Rank of $\begin{pmatrix}1&2&1\\2&4&2\\3&6&3\end{pmatrix}$.
  \item Rank of $\begin{pmatrix}1&0&1\\0&1&1\\1&1&0\end{pmatrix}$.
  \item Rank of $\begin{pmatrix}1&2&3&4\\0&1&2&3\\0&0&0&1\end{pmatrix}$.
  \item Invert $\begin{pmatrix}2&1\\1&1\end{pmatrix}$ by Gauss--Jordan.
  \item Invert $\begin{pmatrix}1&2&3\\0&1&4\\0&0&1\end{pmatrix}$ by
        Gauss--Jordan.
  \item Show $\begin{pmatrix}1&2&3\\4&5&6\\7&8&9\end{pmatrix}$ is singular
        by row reduction.
  \item For which $t$ is $\begin{pmatrix}1&2\\2&t\end{pmatrix}$ singular,
        and what does that mean for the system with that coefficient
        matrix?
  \item Show that $x+y+z=1$, \;$2x+2y+tz=2$ has infinitely many solutions
        for \emph{every} $t$, and say what changes at $t=2$.
  \item Compute $\det\begin{pmatrix}0&1&2\\1&0&3\\2&3&0\end{pmatrix}$ by
        row reduction, tracking the swaps.
  \item After row-reducing a $3\times3$ matrix with one swap and no
        scalings, the diagonal is $(1,2,3)$. Give $\det A$.
  \item Explain why a square system $A\mathbf{x}=\mathbf{b}$ has a unique
        solution for \emph{every} $\mathbf{b}$ exactly when
        $\mathrm{rk}(A)=n$.
  \item A system of $3$ equations in $3$ unknowns has a unique solution.
        What can be said about $\det A$, about the rank, and about the
        homogeneous system with the same coefficients?
  \item \textbf{Leontief, three sectors.} Solve
        $(I-A)\mathbf{x}=\mathbf{d}$ for
        $A=\begin{pmatrix}0&0.5&0\\0.5&0&0.5\\0&0.5&0\end{pmatrix}$ and
        $\mathbf{d}=(10,20,10)^{T}$.
  \item \textbf{Portfolio.} Three assets return $2\%$, $5\%$ and $9\%$. A
        total of $100$ is invested, the whole portfolio must return
        $5\%$, and twice as much goes into asset 1 as into asset 3.
        Find the three amounts.
  \item \textbf{Two-good equilibrium with a tax.} Solve
        $q_{1}+2q_{2}=100$, \;$2q_{1}+q_{2}=110$, then say how the solution
        changes if the second right-hand side rises to $120$, and
        interpret.
\end{problems}

\section{Additional Solutions}

\begin{solutions}
  \item Swap $R_{1}\leftrightarrow R_{2}$, then $R_{3}\to R_{3}-2R_{1}$
        gives a zero row:
        $\begin{pmatrix}1&2&3\\0&1&2\\0&0&0\end{pmatrix}$. Rank $2$.
  \item Continue with $R_{1}\to R_{1}-2R_{2}$:
        $\begin{pmatrix}1&0&-1\\0&1&2\\0&0&0\end{pmatrix}$.
  \item From the second, $x=y-1$; then $3y-3+2y=7$, $y=2$, $x=1$.
  \item Eliminating: $R_{2}-R_{1}$ gives $-2y+z=3$; $R_{3}-2R_{1}$ gives
        $-y-3z=-3$. From the first of these $z=3+2y$; substituting,
        $-y-9-6y=-3$ so $-7y=6$, $y=-\tfrac67$,
        $z=3-\tfrac{12}{7}=\tfrac97$, and
        $x=1-y-z=1+\tfrac67-\tfrac97=\tfrac47$.
  \item All three equations are multiples of $x+2y-z=1$. One pivot, two
        free variables: $x=1-2s+t$, $y=s$, $z=t$.
  \item Adding all three gives $2(x+y+z)=6$, so $x+y+z=3$. Subtracting each
        equation in turn: $z=2$, $x=1$, $y=0$. Unique: $(1,0,2)$.
  \item $R_{2}-4R_{1}$ gives $-3y-6z=-2$; $R_{3}-7R_{1}$ gives
        $-6y-12z=-3$. Doubling the first of these gives $-6y-12z=-4\neq-3$:
        \textbf{inconsistent}, no solution.
  \item The second equation is twice the first, so one pivot and two free
        variables: $x=-2s-3t$, $y=s$, $z=t$.
  \item Two equations, three unknowns. From the second,
        $x_{1}=\tfrac{3-x_{3}}{2}$; substituting into the first,
        $\tfrac{3-x_{3}}{2}-x_{2}+2x_{3}=1$, so
        $x_{2}=\tfrac12+\tfrac32x_{3}$. With $x_{3}=t$:
        $\bigl(\tfrac32-\tfrac t2,\ \tfrac12+\tfrac{3t}{2},\ t\bigr)$.
  \item Adding gives $2x+2z=8$, so $x=4-z$; subtracting gives $2y=4$, so
        $y=2$. With $z=t$: $(4-t,\,2,\,t)$.
  \item Rows $2$ and $3$ are multiples of row $1$: rank $1$.
  \item $R_{3}-R_{1}$ gives $(0,1,-1)$; then $R_{3}-R_{2}$ gives
        $(0,0,-2)$. Three pivots: rank $3$.
  \item Already in REF with three non-zero rows: rank $3$.
  \item $\left(\begin{array}{cc|cc}2&1&1&0\\1&1&0&1\end{array}\right)$;
        swap, then eliminate to get
        $A^{-1}=\begin{pmatrix}1&-1\\-1&2\end{pmatrix}$,
        which checks against
        $\begin{pmatrix}2&1\\1&1\end{pmatrix}
         \begin{pmatrix}1&-1\\-1&2\end{pmatrix}=I$.
  \item Upper triangular with unit diagonal, so clear upwards:
        $A^{-1}=\begin{pmatrix}1&-2&5\\0&1&-4\\0&0&1\end{pmatrix}$.
  \item $R_{2}-4R_{1}$ and $R_{3}-7R_{1}$ give $(0,-3,-6)$ and
        $(0,-6,-12)$; then $R_{3}-2R_{2}$ gives a zero row. Rank $2<3$, so
        the matrix is singular.
  \item $\det=t-4$, so singular exactly at $t=4$. Then the system either
        has no solution or infinitely many, depending on
        $\mathbf{b}$ --- never exactly one.
  \item $R_{2}\to R_{2}-2R_{1}$ gives $(0\ 0\ t-2\mid 0)$, i.e.
        $(t-2)z=0$. For $t\neq2$ this forces $z=0$, leaving $x+y=1$ with
        one free variable --- infinitely many solutions. For $t=2$ the row
        vanishes entirely, leaving $x+y+z=1$ with \emph{two} free
        variables --- again infinitely many, but a larger family. So the
        count of solutions is infinite throughout; what changes at $t=2$ is
        the \textbf{dimension} of the solution set, from a line to a
        plane, because the second equation stops carrying information.
  \item Swap $R_{1}\leftrightarrow R_{2}$ (sign flips):
        $\begin{pmatrix}1&0&3\\0&1&2\\2&3&0\end{pmatrix}$; then
        $R_{3}-2R_{1}$ gives $(0,3,-6)$ and $R_{3}-3R_{2}$ gives
        $(0,0,-12)$. The triangular product is $1\cdot1\cdot(-12)=-12$,
        and one swap flips the sign: $\det A=12$.
  \item One swap contributes a factor $-1$, so
        $\det A=-(1\cdot2\cdot3)=-6$.
  \item $\mathrm{rk}(A)=n$ means $n$ pivots, so no free variables and no
        zero row on the left --- hence no row of the form
        $(0\cdots0\mid\beta)$ can ever appear, whatever $\mathbf{b}$ is.
        Conversely if $\mathrm{rk}(A)<n$ there is a free variable and the
        solution, when it exists, is never unique.
  \item $\det A\neq0$; $\mathrm{rk}(A)=3$ (full rank); and the homogeneous
        system $A\mathbf{x}=\mathbf{0}$ has only the trivial solution
        $\mathbf{x}=\mathbf{0}$ --- all three statements are equivalent to
        invertibility.
  \item $I-A=\begin{pmatrix}1&-0.5&0\\-0.5&1&-0.5\\0&-0.5&1\end{pmatrix}$.
        Scaling by $2$: $2x_{1}-x_{2}=20$, $-x_{1}+2x_{2}-x_{3}=40$,
        $-x_{2}+2x_{3}=20$. By symmetry try $x_{1}=x_{3}$: then the middle
        equation gives $2x_{2}-2x_{1}=40$, and the first gives
        $x_{2}=2x_{1}-20$; substituting, $4x_{1}-40-2x_{1}=40$, so
        $x_{1}=40$, $x_{3}=40$, $x_{2}=60$. Gross outputs $(40,60,40)$.
        (Check the middle row: $-40+120-40=40$ \checkmark.)
  \item With amounts $a_{1},a_{2},a_{3}$:
        $a_{1}+a_{2}+a_{3}=100$, \;$0.02a_{1}+0.05a_{2}+0.09a_{3}=5$,
        \;$a_{1}=2a_{3}$. Substituting the third into the first two:
        $3a_{3}+a_{2}=100$ and $0.04a_{3}+0.05a_{2}+0.09a_{3}=5$, i.e.
        $0.13a_{3}+0.05a_{2}=5$. From the first, $a_{2}=100-3a_{3}$; then
        $0.13a_{3}+5-0.15a_{3}=5$ gives $-0.02a_{3}=0$, so $a_{3}=0$,
        $a_{1}=0$, $a_{2}=100$. The only portfolio meeting all three
        conditions puts everything into asset 2 --- the constraints leave no
        freedom, which the elimination reveals immediately.
  \item Eliminating: from the second, $q_{2}=110-2q_{1}$; substituting,
        $q_{1}+220-4q_{1}=100$, so $-3q_{1}=-120$, $q_{1}=40$ and
        $q_{2}=30$. With $120$ in place of $110$: $q_{2}=120-2q_{1}$ and
        $q_{1}+240-4q_{1}=100$, so $q_{1}=\tfrac{140}{3}\approx46.7$ and
        $q_{2}=120-\tfrac{280}{3}=\tfrac{80}{3}\approx26.7$. The shock to
        one equation moves \emph{both} quantities, one up and one down:
        in a simultaneous system nothing adjusts in isolation.
\end{solutions}
